\section{Piezoelectric shear-horizontal surface waves}\label{sec:bg}

The thermopiezoelectric coupling of \cref{sec:multi} changes bulk stiffness, but piezoelectric media also support surface waves that have no elastic counterpart. The simplest is the Bleustein--Gulyaev (BG) wave, a shear-horizontal (SH) wave localised at the surface of a hexagonal (6mm) crystal. It is included here because it is the prototype of electromechanical surface localisation and because its closed form allows a clean statement of how the electrical boundary condition enters.

\subsection{Formulation}
Let the half-space $y>0$ be a 6mm crystal with the polar axis along $z$ and motion $u_z(x,y,t)$ polarised along $z$, with wave number $k$ along $x$. The relevant constants are the shear stiffness $c_{44}$, piezoelectric constant $e_{15}$ and permittivity $\varepsilon$, combined in $K^2=e_{15}^2/(c_{44}\varepsilon)$ and the stiffened modulus $\bar c=c_{44}(1+K^2)$ with bulk speed $\bar v=\sqrt{\bar c/\rho}$. In the quasi-static approximation the fields are $u=A\,e^{-\alpha y}$ and $\varphi=(e_{15}/\varepsilon)\,u+B\,e^{-ky}$, with $\alpha=k\sqrt{1-v^2/\bar v^2}$, and in the vacuum $y<0$ the potential is $\varphi=C\,e^{ky}$. The boundary conditions at $y=0$ are a free surface $\sigma_{yz}=c_{44}u_{,y}+e_{15}\varphi_{,y}=0$ together with either of the following:
\begin{itemize}
\item \emph{metallised surface}, $\varphi=0$;
\item \emph{free (unmetallised) surface}, continuity of $\varphi$ and of $D_y$ across the interface with vacuum permittivity $\varepsilon_0$.
\end{itemize}

\subsection{Results}
Eliminating the amplitudes gives the closed forms
\begin{equation}\label{eq:bg}
\frac{\alpha}{k}=\frac{K^2}{1+K^2}\ \ (\text{metallised}),\qquad
\frac{\alpha}{k}=\frac{\varepsilon_0}{\varepsilon+\varepsilon_0}\,\frac{K^2}{1+K^2}\ \ (\text{free}),
\end{equation}
so that the surface wave speed is $v=\bar v\sqrt{1-(\alpha_\ast/k)^2}$ with $\alpha_\ast/k$ from \cref{eq:bg}. These were checked against the numerical roots of the $3\times3$ secular determinant (script \path{scripts/check_bg.py}), which agree to the digits quoted. For $K^2=0.2$, $\bar v=1.0954\,v_t$ (with $v_t=\sqrt{c_{44}/\rho}$), the metallised speed is $1.0801\,v_t$, i.e.\ $1.4\%$ below $\bar v$, while for the free surface $v=1.0954$, $1.0953$ and $1.0916\,v_t$ for $\varepsilon_0/\varepsilon=0.01,\,0.1$ and $1$. For $K^2=0.5$ the metallised value is $1.1547\,v_t$ against $\bar v=1.2247\,v_t$.

Two consequences follow. First, the shift between free and metallised speeds, the usual electromechanical coupling measure of a surface wave, is of order $K^4$ and therefore small when $K^2$ is small. Second, the decay length $1/\alpha$ is very sensitive to the electrical condition: for $K^2=0.2$ it is about $1/k\cdot6$ (roughly one wavelength) when metallised, but about $600/k$ (roughly $100$ wavelengths) for a free surface with $\varepsilon_0/\varepsilon=0.01$. A plate of a few wavelengths thickness therefore does not support the free-surface BG wave as a surface wave in a strict sense.

\subsection{Thermal coupling}
For this SH polarisation with uniform fields in $z$ the shear strains $\epsilon_{xz},\epsilon_{yz}$ do not couple to the (diagonal) thermal-stress tensor of a hexagonal crystal, and the pyroelectric coupling acts along the polar axis $z$, which is perpendicular to the in-plane electric field. The BG wave is therefore not coupled to temperature in this geometry, and thermal effects enter only through the temperature dependence of $c_{44}$, $e_{15}$ and $\varepsilon$, which shifts the speed without introducing a loss. Coupled thermo-piezoelectric surface waves exist for other cuts and polarisations, where the sagittal-plane motion has a dilatational component, and for those the Rayleigh-type determinant of \cref{sec:surface} has to be extended by the electrical conditions.
